\documentclass[10pt,a4paper]{amsart}
\usepackage[margin=19mm]{geometry}
\usepackage{amsmath,amssymb,mathtools}
\usepackage[scr=rsfs,cal=euler]{mathalfa}
\usepackage{xfrac}
\usepackage[unicode,hidelinks,bookmarks=false]{hyperref}
\newcommand{\ie}{i.e.\ }
\newcommand{\cf}{cf.\ }
\newcommand{\resp}{respectively\ }
\newcommand{\Subsection}{\subsection}
\providecommand{\nobreakdash}{\nobreak\hbox{-}}
\setlength{\emergencystretch}{2em}
\multlinegap0pt
\usepackage{bm}
\usepackage{bbm}
\usepackage{dsfont}
\usepackage{xspace}
\usepackage{cancel}
\usepackage{mathtools}
\usepackage[shortlabels]{enumitem}
\usepackage{amsthm}
\makeatletter
\@ifundefined{theorem}{\newtheorem{theorem}{Theorem}}{}
\@ifundefined{lemma}{\newtheorem{lemma}[theorem]{Lemma}}{}
\@ifundefined{proposition}{\newtheorem{proposition}[theorem]{Proposition}}{}
\makeatother
\newcommand{\ba}{\mathbf{a}}
\newcommand{\bb}{\mathbf{b}}
\title[Hausdorff dimension of double base expansions and binary shi]{Hausdorff dimension of double base expansions and binary shifts with a hole}
\author{Jian Lu, Wolfgang Steiner, Yuru Zou}
\date{Source: arXiv:2509.04227v2 (CC BY 4.0)}
\begin{document}
\maketitle

\noindent\emph{Source and licence.} Hausdorff dimension of double base expansions and binary shifts with a hole. Version arXiv:2509.04227v2 (CC BY 4.0), distributed under the Creative Commons Attribution 4.0 International licence, as recorded in the arXiv licence metadata for this exact version and shown on the abstract page https://arxiv.org/abs/2509.04227v2. Full rights notice: licenses/arxiv-2509.04227-CC-BY-4.0.txt.\par\smallskip

\noindent\textit{Editorial scope.} Counted unit: the Proposition whose source label is \texttt{p:together} in the article (source0.tex lines 819--822, source label \texttt{p:together}), delivered together with the article's own complete original proof of it (source0.tex lines 824--893, one \texttt{proof} environment of 70 lines). No mathematics is written, repaired or re-derived: statement and proof are the article's own text line for line. Required context retained uncounted: l:ab (lines 584--594); l:phi (lines 632--639); l:phik (lines 677--684); l:sclose (lines 700--703); p:nontogether (lines 748--751); e:anj (lines 780--785). Every definition, display and label of those lines is retained unchanged. Omitted from this package and not redistributed: 1--583, 595--631, 640--676, 685--699, 704--747, 752--779, 786--818, 823--823, 894--1286. Nothing inside a retained excerpt is omitted or altered: every retained body is delivered exactly as the article's lines are written, so no omission receipt of any kind accompanies this package. Citations: the retained excerpts carry no citation command at all, so no bibliography entry of the article is transcribed and no reference is redistributed. Reference adaptation: none was needed and none was made; every reference inside the delivered unit resolves inside it, so no unresolved reference or citation is produced. Printed slips kept literal: no printed word, symbol, hypothesis or number of the counted statement or of its proof was altered. Layout and class: the article's own class is \texttt{amsart} with packages including tikz, hyperref, amsbsy, cite, enumerate; the wrapper is the lane's pinned \texttt{amsart} editorial wrapper at 10pt on A4, which restates the theorem-like environments the retained text needs and adds the article's own macro definitions that the retained text needs (\texttt{\textbackslash ba}, \texttt{\textbackslash bb}), with extra wrapper packages (bm, bbm, dsfont, xspace, cancel, mathtools, enumitem). The delivered numbering is the unit's own. Assets: no retained span contains a figure environment, a graphics-inclusion command, a PDF-page inclusion, a table or a TikZ picture; no image file is copied into this package, the package declares no asset dependency and no image file is redistributed with it. Licence: version arXiv:2509.04227v2 of this article is distributed under the Creative Commons Attribution 4.0 International licence, as recorded in the arXiv licence metadata for this exact version and shown on the abstract page https://arxiv.org/abs/2509.04227v2. A full rights notice is delivered at licenses/arxiv-2509.04227-CC-BY-4.0.txt. Characters: every retained excerpt is pure ASCII, so the delivered engine, XeLaTeX with its default Latin Modern text font, reports no missing character.
\begin{lemma} \label{l:ab}
Let $\ba = a_1a_2\cdots \preceq\bb = b_1b_2\cdots$.
Then we have the following. 
\begin{enumerate}[\upshape (i)]
\itemsep.5ex
\item \label{i:ab1}
 If $a_{n+1}a_{n+2}\cdots \in{]\ba,\bb[}$ for some $n \ge 1$, then $\Omega_{\ba,\bb} = \Omega_{(a_1\cdots a_n)^\infty,\bb}$. 
\item \label{i:ab2}
If $b_{n+1}b_{n+2}\cdots \in {]\ba,\bb[}$ for some $n \ge 1$, then $\Omega_{\ba,\bb} = \Omega_{\ba,(b_1\cdots b_n)^\infty}$. 
\end{enumerate}
\end{lemma}

\begin{lemma} \label{l:phi}
Let $\varphi$ be a substitution such that $\varphi(0) \in 0\{0,1\}^*$ and $\varphi(1) \in 1\{0,1\}^*$, $\ba \in 0\{0,1\}^\infty$, $\bb \in 1\{0,1\}^\infty$.
Then
\[
0 \le d_{q_0,q_1}(\varphi(\ba),\varphi(\bb)) - d_{q_0,q_1}(\varphi(0^\infty),\varphi(1^\infty)) \le \dim_H \pi_{q_0,q_1}\big(\varphi(\Omega_{\ba,\bb})\big).
\]
In particular, $d_{q_0,q_1}(\varphi(\ba),\varphi(\bb)) = d_{q_0,q_1}(\varphi(0^\infty),\varphi(1^\infty))$ when $h(\Omega_{\ba,\bb}) = 0$.
\end{lemma}

\begin{lemma} \label{l:phik}
Let $\ba \in 0\{0,1\}^\infty$, $\bb \in 1\{0,1\}^\infty$. 
If there exists a sequence of substitutions $(\varphi_k)_{k\ge1}$ such that the length of  $\varphi_k(01)$ is unbounded and
\begin{equation*} \label{e:phii}
\varphi_k(0^\infty) \prec \ba \prec \varphi_k(01^\infty),\ \varphi_k(10^\infty) \prec \bb \prec \varphi_k(1^\infty) \quad \mbox{for all}\ k \ge 0,
\end{equation*}
then $d_{q_0,q_1}$ is continuous at $(\ba,\bb)$.
\end{lemma}

\begin{lemma} \label{l:sclose}
Let $q_0, q_1 > 1$.
The restriction of the map $d_{q_0,q_1}$ to~$W$ is continuous.
\end{lemma}

\begin{proposition} \label{p:nontogether}
Let $(\ba,\bb) = (a_1a_2\cdots, b_1b_2\cdots) \in W$, with $a_{m+1} a_{m+2} \cdots \ne \bb$ and $b_{m+1} b_{m+2}\cdots \ne \ba$ for all $m \ge 1$.
Then $d_{q_0,q_1}$ is continuous at $(\ba,\bb)$. 
\end{proposition}

\begin{equation} \label{e:anj}
\begin{aligned}
a_{i+1} \cdots a_{n_{j+1}+1} & = b_{i-n_j+1} \cdots b_{n_{j+1}-n_j} 1 \succ b_{i-n_j+1} \cdots b_{n_{j+1}-n_j+1} \\
& \succeq b_1 \cdots b_{n_{j+1}-i+1} = a_{n_j+1} \cdots a_{n_j+n_{j+1}-i+1},
\end{aligned}
\end{equation}

\begin{proposition} \label{p:together}
Let $(\ba,\bb) = (a_1a_2\cdots, b_1b_2\cdots) \in W$ with $a_{m+1} a_{m+2} \cdots = \bb$ or $b_{m+1} b_{m+2} \cdots = \ba$ for some $m \ge 1$.
Then $d_{q_0,q_1}$ is continuous at $(\ba,\bb)$. 
\end{proposition}

\begin{proof}
Assume that $b_{m+1} b_{m+2} \cdots = \ba$, the other case being symmetric.
Moreover, we assume w.l.o.g.\ that $m$ is minimal, i.e., $b_{i+1} b_{i+2} \cdots \ne \ba$ for all $i < m$.

If $a_{i+1}a_{i+2}\cdots = \bb$ for some $i \ge 1$, then $\ba = \varphi(01)^\infty$, $\bb = \varphi(10)^\infty$ with $\varphi(0) = a_1\cdots a_i$, $\varphi(1) = b_1\cdots b_m$.
Since $h(\Omega_{0110^\infty,1001^\infty}) = h(\Omega_{(01)^\infty,(10)^\infty}) = 0$ and $\ba \prec \varphi(0110^\infty)$, $\bb \succ \varphi(1001^\infty)$, $d_{q_0,q_1}$ is constant in a neighbourhood of $(\ba,\bb)$ by Lemma~\ref{l:phi}, thus continuous at $(\ba,\bb)$. 

Assume now that $a_{i+1}a_{i+2}\cdots \ne \bb$ for all $i \ge 1$.
Then we can define $n_k$ as in the proof of Proposition~\ref{p:nontogether}.
For $k \ge 0$, define the substitution~$\varphi'_k$ by
\[
\varphi'_k(0) = a_1 \cdots a_{n_k}, \quad \varphi'_k(1) = b_1 \cdots b_m.
\]
We claim that, for all $k \ge 1$,
\begin{equation} \label{e:phi'}
\big(\varphi'_k(01^\infty), \varphi'_k(1^\infty)\big) \in W.
\end{equation}

To prove the claim, note first that the only sequence of $\Omega_{\ba,\bb}$ starting with $\varphi'_k(1)0$ is~$\bb$.
Since $a_{n+1}a_{n+2}\cdots \ne \bb$ for all $n \ge 1$ and $\ba \in \Omega_{\ba,\bb}$, this implies that $\ba \in \Omega_{\ba,\varphi'_k(1^\infty)}$, thus 
\begin{equation} \label{e:ai}
a_1 \cdots a_i \varphi'_k(1^\infty) \preceq \ba \quad \mbox{for all $i \ge 1$ such that $a_{i+1} = 1$.}
\end{equation}
In particular, we have
\[
\varphi'_k(01^\infty) \preceq \ba
\]
and thus $n_{k+1}{-}n_k \le m$ for all $k \ge 0$.

We strengthen now \eqref{e:ai} to 
\begin{equation} \label{e:ai2}
a_1 \cdots a_i \varphi'_k(1^\infty) \preceq \varphi'_k(01^\infty) \quad \mbox{for all $1 \le i \le n_k$ such that $a_{i+1} = 1$,}
\end{equation}
and, equivalently, 
\begin{equation} \label{e:ai3}
a_{i+1} \cdots a_{n_k} \varphi'_k(1^\infty) \succeq \varphi'_k(1^\infty) \quad \mbox{for all $1 \le i < n_k$ such that $a_{i+1} = 1$.}
\end{equation}
Let $n_j \le i < n_{j+1}$, $0 \le j < k$, $a_{i+1} = 1$.
Then $a_{i+1} \cdots a_{n_{j+1}+1} \succ b_1 \cdots b_{n_{j+1}-i+1}$ by~\eqref{e:anj}, and $n_{j+1}{-}i \le n_{j+1}{-}n_j \le m$ implies that $a_{i+1} \cdots a_{n_{j+1}} \varphi'_k(1^\infty) \succeq \varphi'_k(1^\infty)$.
Recursively from $j = k{-}1$ to $j = 0$, we obtain that
\[
\begin{aligned}
a_{i+1} \cdots a_{n_k} \varphi'_k(1^\infty) & = a_{i+1} \cdots a_{n_{j+1}} a_{n_{j+1}+1} \cdots a_{n_k} \varphi'_k(1^\infty) \\
& \succeq a_{i+1} \cdots a_{n_{j+1}} \varphi'_k(1^\infty) \succeq \varphi'_k(1^\infty).
\end{aligned}
\]
This proves \eqref{e:ai3} and thus \eqref{e:ai2}.

To show that $\varphi'_k(1^\infty) \in \Omega_{\varphi'_k(01^\infty),\varphi'_k(1^\infty)}$, let $1 \le i < m$.
If $b_{i+1} = 1$, then $b_{i+1} \cdots b_m b_1 \succ b_{i+1} \cdots b_{m+1} \succeq b_1 \cdots b_{m-i+1}$, thus $b_{i+1} \cdots b_m \varphi'_k(1^\infty) \succ \varphi'_k(1^\infty)$.
If $b_{i+1} = 0$, then $b_{i+1} \cdots b_{m+1} \preceq a_1 \cdots a_{m-i+1}$, with $a_{m-i+1} = 1$ in case $b_{i+1} \cdots b_m = a_1 \cdots a_{m-i}$ because $a_{m-i+1} = 0$ would imply $\ba \preceq a_1 \cdots a_{m-i} \ba = b_{i+1} b_{i+2} \cdots \preceq \ba$, thus $b_{i+1} b_{i+2} \cdots = \ba$, contradicting the minimality of~$m$.
By \eqref{e:ai2}, which we can use because $m < m{+}n_{k-1} \le n_k$, we obtain that $b_{i+1} \cdots b_m \varphi'_k(1^\infty) \preceq \varphi'_k(01^\infty)$, thus $\varphi'_k(1^\infty) \in \Omega_{\varphi'_k(01^\infty),\varphi'_k(1^\infty)}$.
To prove \eqref{e:phi'}, it remains to consider $a_{i+1} \cdots a_{n_k} \varphi'_k(1^\infty)$ for $1 \le i < n_k$.
If $a_{i+1} = 0$, then $a_{i+1} \cdots a_{n_k+1} \preceq a_1 \cdots a_{n_k-i+1}$,  
$a_{n_k+1} = 1$, and \eqref{e:ai2} give that $a_{i+1} \cdots a_{n_k} \varphi'_k(1^\infty) \preceq \varphi'_k(01^\infty)$.
By \eqref{e:ai3}, we obtain $\varphi'_k(01^\infty) \in \Omega_{\varphi'_k(01^\infty),\varphi'_k(1^\infty)}$, i.e., \eqref{e:phi'} holds.

Consider now $\ba' \succeq \varphi'_k(01^\infty)$, $\bb' \preceq \varphi'_k(1^\infty)$.
Then $\ell(\ba',\bb') \succeq \varphi'_k(01^\infty)$ and $r(\ba',\bb') \preceq \varphi'_k(1^\infty)$.
In particular, $\ell(\ba',\bb')$ is close to~$\ba'$ for $\ba'$ close to the right of $\varphi'_k(01^\infty)$.
Since $b_{i+1} \cdots b_m \ba \succ \bb$ for all $1 \le i < m$ such that $b_{i+1} = 1$, we have $b_{i+1} \cdots b_m \ell(\ba',\bb') \succeq \bb'$ for all such~$i$ when $\ell(\ba',\bb')$ is close to~$\ba$ and $\bb'$ is close to~$\bb$.
Since $b_{i+1} \cdots b_m \ell(\ba',\bb') \preceq b_{i+1} \cdots b_m \varphi'_k(1^\infty) \preceq \varphi'_k(01^\infty)$ for $b_{i+1} = 0$, this implies that $b_{i+1} \cdots b_m \ell(\ba',\bb') \notin {]\ba',\bb'[}$ for all $1 \le i < m$.
If $\bb' \preceq \varphi'_k(1) \ell(\ba',\bb')$, then we obtain that $\varphi'_k(1) \ell(\ba',\bb') \in \Omega_{\ba',\bb'}$ and thus $r(\ba',\bb') \preceq \varphi'_k(1) \ell(\ba',\bb')$, hence $r(\ba',\bb')$ is close to~$\bb'$, and $d_{q_0,q_1}(\ba',\bb')$ is close to $d_{q_0,q_1}(\ba,\bb)$ by Lemma~\ref{l:sclose}.
If $\bb' \succ \varphi'_k(1) \ell(\ba',\bb')$, then Lemma~\ref{l:ab} gives that $r(\ba',\bb') = \varphi'_k(1^\infty)$.
Since $\Omega_{\ba,\bb} \setminus \Omega_{\ba,\varphi'_k(1^\infty)}$ contains only sequences ending with~$\bb$, we have $d_{q_0,q_1}(\ba,\bb) = d_{q_0,q_1}(\ba,\varphi'_k(1^\infty))$, and again $d_{q_0,q_1}(\ba',\bb')$ is close to $d_{q_0,q_1}(\ba,\bb)$ by Lemma~\ref{l:sclose}.
This proves that, for large~$k$, $d_{q_0,q_1}(\ba',\bb')$ is close to $d_{q_0,q_1}(\ba,\bb)$ for all $\ba'$ close to the right of $\varphi'_k(01^\infty)$ and close to~$\bb$.
If $\varphi'_k(01^\infty) \ne \ba$ for all~$k$, i.e., $\varphi'_k(01^\infty) \nearrow \ba$, then this proves the continuity of $d_{q_0,q_1}$ at $(\ba,\bb)$.
If $\varphi'_k(01^\infty) = \ba$ for some $k \ge 1$, then this holds for all sufficiently large~$k$, and the proof of Lemma~\ref{l:phik} with the substitutions~$\varphi'_k$ shows that $d_{q_0,q_1}(\ba',\bb')$ is close to $d_{q_0,q_1}(\ba,\bb)$ for all $\ba'$ close to the left of~$\ba$ and close to~$\bb$.
Therefore, $d_{q_0,q_1}$ is continuous at $(\ba,\bb)$ also in this case.
\end{proof}

\end{document}
